\documentclass[letterpaper]{article}
\usepackage{amsmath,amssymb}
\usepackage[letterpaper,margin=1in]{geometry}
\setlength{\parindent}{0pt}
\newcommand{\noteheader}[1]{\par\bigskip\noindent\textbf{#1}\par\medskip}

\begin{document}

\textbf{CS 2050 --- Notes 09/30}

\noteheader{Union of sets}

For sets $A$ and $B$, the \textbf{union} $A\cup B$ contains the elements
that are in $A$, in $B$, or in both:
\[
A\cup B=\{x\mid x\in A\lor x\in B\}.
\]

\noteheader{Intersection of sets}

The \textbf{intersection} $A\cap B$ contains the elements that are in
both $A$ and $B$:
\[
A\cap B=\{x\mid x\in A\land x\in B\}.
\]

\noteheader{Disjoint sets}

Two sets are \textbf{disjoint} if their intersection is empty:
\[
A\cap B=\varnothing.
\]

\noteheader{Difference of sets}

The \textbf{difference} $A-B$ contains the elements in $A$ that are not
in $B$. This is the complement of $B$ relative to $A$:
\[
A-B=\{x\mid x\in A\land x\notin B\}.
\]

\clearpage

\noteheader{Complement of a set}

Let $U$ be the universal set and $A\subseteq U$. The \textbf{complement}
of $A$ relative to $U$ contains the elements of $U$ that are not in $A$:
\[
A^c=U-A=\{x\in U\mid x\notin A\}.
\]
The notations $\overline A$ and $A^c$ both denote this complement.

\end{document}
